\hspace{3mm}

\centerline{\bf \Huge Домашнее задание.}
\centerline{\bf \Huge Принцип \href{https://vk.com/video-91031095_456245511}{Дирихлееее!}}

\begin{flushright}
    \scriptsize{Чтобы найти решение, найти которое невозможно, \\
                нужно допускать ошибки, иначе все бы знали, как оно решается\\
                Барбосса мира ЭлМШ}
\end{flushright}

\problem{1!} В шиншиллах обучаются 27 учеников. Докажите, что по крайней мере двое учеников родились в один и тот же час. (Например, АА родился в 14 часов дня и РВ родился в 2 часа ночи. АА и РВ родились в разное время)

\textit{Комментарий от РВ:} я родился в 15:20

\problem{2!} В первой задаче Олимпиады ЭлМШ было написано 40 слов. Докажите, что хотя бы два слова заканчивались на одинаковую букву.

\problem{3} В ЭлМШ обучаются 50 учеников. Докажите, что хотя бы 13 учеников родились в одно время года.

\problem{4} В олимпиаде ЭлМШ 5 задач стоимость каждой 7 баллов. Докажите, что хотя бы два ученика ЭлМШ набрали одинаковое количество баллов, если от ЭлМШ участвовало 37 учеников.